\documentclass[10pt]{article}
\usepackage[margin=0.85in]{geometry}
\usepackage{amsmath,amssymb,booktabs,tabularx,array,hyperref}
\newcolumntype{Y}{>{\raggedright\arraybackslash}X}
\title{Latent Cognitive Blocks:\\Parallel Causal Span Control for Just-in-Time Coprocessor Invocation}
\author{Working Draft}\date{August 2026}
\begin{document}\maketitle

\begin{abstract}
Explicit cognitive blocks serialize control through language tokens. We
pre-register a parallel causal channel that marks the start, coarse type, and
end of a cognitive-assistance span while ordinary autoregressive generation
continues. The stable controller vocabulary is \textsc{none}, \textsc{compute},
\textsc{retrieve}, \textsc{verify}, and \textsc{help}; a second-stage runtime
resolves deployment-specific capabilities. This is a protocol paper, not a
positive latent-control result. Paper 2's engine-specific learned-routing gate
remains \texttt{no\_go}. The independent latent experiment may proceed only as
a matched comparison against generic paired tags, fenced blocks, label-only
intent, transparent CPU triggers, and oracle spans. We separate causal span
control from payload normalization, deterministic execution, result
integration, and optional PRA materialization.
\end{abstract}

\section{Status and Motivation}
Paper 2 shows that deterministic parsing and execution are reliable once a
capability is known, but learned engine identity is unsafe under held-out
controls and registry growth. Its machine-readable Paper 3 gate is
\texttt{no\_go}; that decision is not overwritten here. The new question is
narrower and registry-independent: can a frozen Transformer expose when a span
needs cognitive assistance, what coarse assistance family it needs, and when
the span is complete, without serializing control syntax into generated text?

The hypothesis concerns transport across the neural--runtime boundary, not a
larger engine classifier. Conventional explicit tools remain the universal
fallback. Runtime-copy remains the default for authoritative exact outputs.

\section{Factorized Architecture}
At token $t$ and selected layer $l$, a frozen representation $h_t^{(l)}$ feeds a
small causal controller. Outside a block,
\[
p_{\rm start}^{(l)}=\operatorname{softmax}(W_{\rm start}^{(l)}h_t^{(l)})
\]
ranges over \textsc{none}, \textsc{compute}, \textsc{retrieve},
\textsc{verify}, and \textsc{help}. Inside a block, a separate head predicts
\textsc{continue} or \textsc{end}. No future token may inform a headline
decision.

The two routing stages are
\[
R_1:\text{generation}\rightarrow
\{\textsc{none},\textsc{compute},\textsc{retrieve},\textsc{verify},\textsc{help}\},
\]
\[
R_2:(\text{intent},\text{span},\text{registry},\text{policy},\text{state})
\rightarrow\text{capability or shortlist}.
\]
$R_1$ never predicts calculator, graph, Datalog, DuckDB, or another registry
identity. $R_2$ may use deterministic rules, token-aware BM25, a CPU classifier,
capability retrieval, optional PRA, or an explicit tool fallback.

\subsection{Constrained state machine}
The preferred minimal machine has two states:
\begin{itemize}
\item \textbf{OUTSIDE}: emit \textsc{none} or
\textsc{start}$_{\{\rm compute,retrieve,verify,help\}}$;
\item \textbf{INSIDE(type)}: emit \textsc{continue} or \textsc{end}.
\end{itemize}
Nesting and type switching are invalid. On \textsc{end}, generation pauses;
the runtime extracts the naturally generated span, resolves and executes a
capability, records provenance-bearing typed state, materializes a bounded
result, and resumes OUTSIDE. BIO/BIES labels are retained as a matched
alternative, not assumed superior.

\section{Payload and Completion Ladders}
The payload ladder starts with the marked ordinary-language span. If it is not
directly executable, the runtime tries a deterministic parser, then a selected
PRA/ICL capability skill, then short capability-specific normalization. A
learned structured-argument head is justified only after those fail. This keeps
span control distinct from normalization quality.

For \textsc{end}, compare neural prediction, deterministic parser completeness,
their conservative AND, and a safe OR that fires only when either signal is
high-confidence and policy permits. Arithmetic, date, and units may rely more
on deterministic closure than \textsc{retrieve} or \textsc{help}. Report
premature and late termination separately.

\section{Multi-Layer Control}
Probe approximately $0.25L$, $0.40L$, $0.55L$, $0.70L$, $0.85L$, and $L$.
Final-layer control is only the baseline. The registered aggregation ladder is:
best development-selected layer, fixed mean of layer logits, majority or
threshold vote, learned scalar layer weights, and a tiny MLP over layer logits.
Hidden-state fusion is deferred unless logit aggregation fails. START and END
may select different fixed mixtures, but every choice is frozen on development
data before test evaluation.

\section{Matched Benchmark and Baselines}
Create one new causal-span freeze spanning calculator, date, units, graph,
Datalog, retrieval, verification, ambiguous \textsc{help}, and \textsc{none}.
Controls include quoted tags/fences, discussion of tools without a current need,
ordinary prose, and long generations with multiple opportunities. Split by
paraphrase family and namespace, not random surface form. Derive explicit and
latent labels from the same examples.

\begin{table}[t]
\centering\small
\caption{All interfaces use the same examples and runtime capabilities.}
\begin{tabularx}{\linewidth}{@{}p{0.25\linewidth}Y@{}}
\toprule
Condition & Question isolated \\
\midrule
Paired generic tags & Explicit bounded language-bus control \\
Fenced generic block & Familiar delimiter with Markdown collision risk \\
Label-only generic intent & Minimum generated control tokens; active task is payload \\
Transparent CPU trigger & Non-neural causal control baseline \\
Latent span controller & Parallel START/type/END control bus \\
Oracle span and type & Upper bound before R2/normalization/execution \\
\bottomrule
\end{tabularx}
\end{table}

\subsection{Cross-paper public control registry}
To avoid inventing a new task distribution, we freeze 280 redistribution-safe
references from retained Papers 1.5, 2, and 2.5 selections: 40 each from GSM8K,
BIG-bench Date Understanding, BIG-bench Unit Conversion, balanced ProofWriter,
CLUTRR, CRAG, and TAT-QA. The artifact stores only source identifiers, benchmark
IDs, content hashes, deterministic selection keys, coarse assistance type, and
span-audit status. TAT-QA arithmetic/count rows record \textsc{compute} as a
secondary action after primary \textsc{retrieve}.

This registry is not yet a latent-control benchmark. Every row is marked
\texttt{pending\_audit}; a bounded FRAMES subset is also absent. Before any
headline run, auditors must add earliest useful START, latest safe START,
canonical END, payload span, assistance type, and first wrong or unsupported
token where applicable. The manifest therefore records
\texttt{headline\_ready: false} and blocks result generation rather than
synthesizing span labels from answers.

The required public matrix is four generic tools, paired tags, fenced blocks,
label-only intent, CPU trigger, final-layer latent, multi-layer latent,
tools-plus-watchdog, and oracle timing/type, all with identical R2 and backends.
The oracle-timed Paper 2 and 2.5 gateway audits establish only transport/schema
parity; they do not substitute for voluntary or latent public runs.

\section{Training and Registry Invariance}
The ladder is P0 frozen linear probes, P1 frozen multi-layer ensembles, P2 a
small control-only LoRA if necessary, and P3 optional multitask LM/control
training only after a positive gate. Train with calculator, date, and units in
the registry; test $R_1$ unchanged after adding graph and Datalog, then later
SMT, algebra, or retrieval backends. Score $R_1$ intent and $R_2$ capability
selection separately. Adding an implementation must not silently alter test
prompts or retrain the generic controller.

Result assimilation is a separate factorial axis: token versus latent control,
crossed where feasible with base versus integration LoRA. A control gain must
not be inferred from runtime-copy, and an integration gain must not be called a
better interrupt.

\section{Metrics and Gate}
Report START/END precision, recall, and F1; exact span and token IoU; coarse type
accuracy; NONE false-activation rate; premature/late END; payload parseability;
R2 capability success; runtime and final exactness; controller parameters,
FLOPs, and latency; explicit control tokens avoided; earliest reliable layer;
cross-layer agreement and calibration; and registry-expansion deltas. Preserve
raw causal traces and multi-seed uncertainty for headline comparisons.

Latent control advances only if it improves reliability or matched cost over
generic explicit blocks and transparent CPU triggers without unsafe FAR or
brittle END behavior. It is falsified as a needed mechanism if explicit generic
blocks are equally reliable and cheap, aggregation adds no value, coupling is
model-specific, or CPU triggers dominate. Early exit is exploratory only after
strong multi-layer evidence and receives no headline claim in this version.

\section{Scope, PRA, and Execution Order}
There are no nested latent blocks, general planner, learned engine registry,
joint base-model training, backtracking contribution, or silent promotion of
uncertain language to fact. PRA is optional: intent may retrieve and materialize
one to $k$ capability skills, but the generic controller must work without it.

The execution order is dataset freeze; shared explicit/latent labels;
final-layer probe; layer sweep; aggregation; BIO versus START/END;
parser-assisted END; matched syntax and CPU comparisons; registry expansion;
control LoRA only if needed; separate integration LoRA; and only then an
early-exit study. Until retained artifacts complete these steps, this manuscript
states hypotheses and protocol rather than empirical superiority.

The public registry can be reproduced from retained selections with:
\begin{verbatim}
python -m ccpu paper3 freeze-public \
  --compute-selection artifacts/paper2/public_benchmarks/data_v1/selection.jsonl \
  --crag-selection artifacts/paper1_5/public_benchmarks/crag_v1/data/selection.jsonl \
  --tatqa-selection artifacts/paper2_5/public_benchmarks/tatqa_v1/selection.jsonl \
  --per-benchmark 40 \
  --output-dir artifacts/paper3/public_control_v1
\end{verbatim}

\begin{thebibliography}{9}
\bibitem{toolformer} T. Schick et al. Toolformer: Language Models Can Teach
Themselves to Use Tools. \emph{NeurIPS}, 2023.
\bibitem{selfrag} A. Asai et al. Self-RAG: Learning to Retrieve, Generate, and
Critique through Self-Reflection. \emph{ICLR}, 2024.
\bibitem{logiclm} L. Pan et al. Logic-LM: Empowering Large Language Models with
Symbolic Solvers for Faithful Logical Reasoning. \emph{EMNLP}, 2023.
\end{thebibliography}
\end{document}
